\section{A positive Green-strip potential and its prime baseline}
\label{sec:green-strip}

A second metric measures nonreal zeros directly and is invariant under
scalar normalization. Throughout this section \(0\le\lambda\le1/2\),
so every canonical zero satisfies \(|\Im z|\le1\), including at time
zero. We retain all multiplicities and the complete opposite-paired
product and counts used above. The complete Green identity and its
arithmetic reduction have distinct internal assessments
\cite{GreenStripProof,LocalizedMetricsAssessment,MetricComparison,PrimeBaselineAssessment}.

\subsection{The exact strip identity and finite-factor proof}
Put
\begin{equation}\label{eq:green-kernels}
 a=\pi/8,\qquad K(t)=\frac1{16\cosh(at)},\qquad
 G(d,y)=\operatorname{atanh}\frac{\sin(ay)}{\cosh(ad)}.
\end{equation}
Each boundary kernel has mass \(\int_\R K=1/2\). Define
\begin{align}\label{eq:green-J}
 J_\lambda(X)={}&\int_\R K(x-X)
       \{\log|H_\lambda(x)|+\log|H_\lambda(x+8\ii)|\}\dd x
       -\log|H_\lambda(X+4\ii)|.
\end{align}
\begin{theorem}[Complete positive strip potential]\label{thm:green-identity}
Both boundary integrals in \eqref{eq:green-J} are absolutely integrable,
including at real multiple zeros, and for every real \(X\),
\begin{equation}\label{eq:green-identity}
 J_\lambda(X)=\sum_{j:y_j(\lambda)>0}
                  G(X-x_j(\lambda),y_j(\lambda)).
\end{equation}
This complete upper-zero sum converges absolutely. Every summand is
positive. Real boundary zeros give no atom, and there is no additional
outer or harmonic term. Multiplication of \(H_\lambda\) by any nonzero
complex scalar leaves \(J_\lambda\) unchanged.
\end{theorem}

We first prove the finite-factor identity and then pay its complete
limit. On \(0<\Im z<B\), with \(-\Delta G_B=2\pi\delta\), the strip
Green function is
\begin{equation}\label{eq:green-general}
 G_B(z,\zeta)=\frac12\log
 \frac{\cosh(\pi(\Re z-\Re\zeta)/B)-\cos(\pi(\Im z+\Im\zeta)/B)}
      {\cosh(\pi(\Re z-\Re\zeta)/B)-\cos(\pi(\Im z-\Im\zeta)/B)}.
\end{equation}
The map \(z\mapsto\ee^{\pi z/B}\) to the upper half-plane gives
this reflected-distance ratio, its zero boundary values and singularity
\(-\log|z-\zeta|\). Boundary normal derivatives divided by \(2\pi\)
give the Poisson kernels. At \(B=8\), observation \(X+4\ii\), both
are \(K(x-X)\), and the Green value is \eqref{eq:green-kernels}.

For an interior zero the harmonic correction is
\(\log|z-\zeta|+G_B(z,\zeta)\); its singularities cancel and its
boundary values are those of the factor logarithm. Its Poisson
representation follows by truncation with vertical sides: their harmonic
measures decay exponentially, while the factor logarithm grows only
logarithmically. Thus the boundary average minus the interior logarithm
is \(+G_B\). Exterior zeros give zero. A boundary zero is treated by its
locally integrable logarithm or an exterior limit and gives no half-weight
atom, also at finite multiple order.

The inherited exact even product is
\begin{equation}\label{eq:green-product}
 H_\lambda(z)=H_\lambda(0)
      \prod_{\zeta\in\mathcal R_\lambda}(1-z^2/\zeta^2),
\end{equation}
where \(\mathcal R_\lambda\) has one representative of each opposite
pair, with multiplicity. No surviving exponential factor is omitted.
For each finite complete opposite-paired product, apply the factor
identity to \(z-\zeta\), \(z+\zeta\) and its constants. The constants
cancel since the boundary masses add to one. This proves the finite
version of \eqref{eq:green-identity}.

\subsection{Full-product logarithmic integrability and boundary tails}
Use a common full-count constant \(C\) such that
\begin{equation}\label{eq:green-full-count}
 N_\lambda(r)\le Cr\log(r+2),\qquad r\ge1.
\end{equation}
The complete source and positivity at the origin give finite constants
\[
 m_0=\inf_\lambda H_\lambda(0)>0,\qquad
 D_0\ge\sup_{\lambda,|z|\le1}|H_\lambda'(z)|,\quad D_0>0.
\]
Set
\begin{equation}\label{eq:green-log-constants}
 r_*=\min\{1,m_0/(2D_0)\},\quad C_0=\sup_\lambda|\log H_\lambda(0)|,
 \quad C_*=C_0+20C(1+|\log r_*|).
\end{equation}
The segment integral of \(H'\) gives \(|H_\lambda(z)|\ge m_0/2\)
on \(|z|\le r_*\), so every zero lies outside this fixed disk.

For every real \(q\), \(0\le b\le8\), and \(Q=|q|+12\),
\begin{equation}\label{eq:green-unit-log}
 \int_q^{q+1}|\log|H_\lambda(x+\ii b)||\dd x
                  \le C_*Q\log^2(3Q).
\end{equation}
To prove the stronger assertion needed for the product limit, bound the
sum of absolute factor logarithms. On the interval \(|x+\ii b|\le Q\).
For \(|\zeta|\le2Q\), split a factor logarithm into two distance logs
and \(-2\log|\zeta|\). Each positive distance log is at most
\(\log(3Q)\), while
\[
 \int_q^{q+1}\log^-|x+\ii b-\zeta|\dd x\le2,
\]
because distance is at least \(|x-\Re\zeta|\), and
\(\int_{-1}^1-\log|t|\dd t=2\). This includes repeated boundary
zeros. Each near representative therefore costs at most
\(4\log(3Q)+4+2|\log r_*|\). There are at most
\(CQ\log(2Q+2)\) such representatives, using half the full
origin-centered count at radius \(2Q\).
For far representatives,
\[
 \left|\log\left|1-\frac{z^2}{\zeta^2}\right|\right|
         \le\frac43\frac{Q^2}{|\zeta|^2}.
\]
Partial summation of \eqref{eq:green-full-count}, again halved for
representatives, bounds their total by
\[
 \frac23CQ\{\log(2Q+2)+1\}.
\]
These costs imply
\eqref{eq:green-unit-log} with \eqref{eq:green-log-constants}.
They also prove absolute local \(L^1\) convergence of the factor
logarithms; the polynomial unit-interval bound is summable against
the exponential Poisson weight. At \(X+4\ii\), the product converges
away from every zero by the common zero strip. Passing from the finite
product identity to the complete one is now justified. This proves
Theorem~\ref{thm:green-identity} with no unaccounted outer term.
Scalar invariance follows because the added \(\log|c|\) has total
boundary mass one and is subtracted once at the observation point.

There is an explicit uniform two-boundary tail. For \(R\ge0\), set
\(n=\lfloor R\rfloor\), \(q_0=\ee^{-a}\), \(S=|X|+13+n\). Then
\begin{align}\label{eq:green-boundary-tail}
 &\int_{|x-X|>R}K(x-X)
   \{|\log|H_\lambda(x)||+|\log|H_\lambda(x+8\ii)||\}\dd x
 \nonumber\\
 &\quad\le\frac{C_*}{2}\ee^{-an}S\log^2(3S)
                        \frac{1+4q_0+q_0^2}{(1-q_0)^4}.
\end{align}
For shell \(k\le|x-X|<k+1\), use two unit intervals,
\(K\le\ee^{-ak}/8\), and both boundaries. The remaining sum uses
\((S+j)\log^2(3(S+j))\le S\log^2(3S)(1+j)^3\) and
\(\sum_{j\ge0}q_0^j(1+j)^3=(1+4q_0+q_0^2)/(1-q_0)^4\).
The displayed height factor must be retained in joint growing-height
and truncation limits.

\subsection{The complete zero tail and a one-defect margin}
For \(0<y\le1\),
\begin{equation}\label{eq:green-point-tail}
 0<G(d,y)\le K_Gy\ee^{-a|d|},\qquad K_G=2a/\cos^2a.
\end{equation}
Indeed \(\operatorname{atanh}u\le u/(1-u^2)\),
\(\sin(ay)\le ay\), and \(\operatorname{sech}(ad)\le2\ee^{-a|d|}\).
Let \(A\) be the full unit-count constant in
\eqref{eq:coarse-canonical}, and put \(Q_X=3+|X|\). For the same
\(n,q_0\) as above, omitted upper zeros obey
\begin{align}\label{eq:green-zero-tail}
 \sum_{\substack{y_j>0\\|x_j-X|>R}}G(X-x_j,y_j)
 \le2AK_G\ee^{-an}\left\{
 \frac{\log Q_X+\log(n+1)}{1-q_0}
 +\frac{q_0}{(n+1)(1-q_0)^2}\right\}.
\end{align}
Cover both sides by unit shells, use
\(\log(Q_X+k)\le\log Q_X+\log(k+1)\), followed by
\(\log(n+j+1)\le\log(n+1)+j/(n+1)\).
This proves complete convergence and
\(0\le J_\lambda(X)\ll_A\log(3+|X|)\), uniformly in time.
The boundary tail \eqref{eq:green-boundary-tail} and the zero tail
\eqref{eq:green-zero-tail} pay different truncations.

If an upper zero has \(y_j\ge d>0\), \(|x_j-X|\le r\), it forces
\begin{equation}\label{eq:green-defect-margin}
 J_\lambda(X)\ge\frac{\sin(a)d}{\cosh(ar)}.
\end{equation}
This uses \(\operatorname{atanh}u\ge u\) and concavity
\(\sin(ay)\ge y\sin a\) on \([0,1]\).
An independently obtained strict upper bound below this margin would
exclude that declared defect class. The acquired coarse bound does not.
An all-real slice gives \(J=0\); any upper zero gives \(J>0\) at every
center. Universal vanishing on the zeta slice is consequently an
RH-type reformulation, stated here without treating checks as evidence
for it.

The complete source supplies
\[
 S_8=\int_\R\ee^{u^2/2+8|u|}\Phi(u)\dd u<\infty,\qquad
 |H_\lambda(x+\ii b)|\le S_8/2\quad(0\le b\le8).
\]
Integrating the safe-line logarithmic derivative gives
\begin{equation}\label{eq:green-safe-lower}
 \log|H_\lambda(X+4\ii)|\ge\log H_\lambda(4\ii)
                              -A|X|\log(2+|X|).
\end{equation}
Neither these separate estimates nor the better count bound yields a
signed boundary/interior comparison below \eqref{eq:green-defect-margin}.
In particular an absolute finite-source error for \(H\) cannot be put
inside its logarithm near real zeros without a relative margin.

Uniform tails on compact center ranges and continuity of finite root
multisets prove continuity of \(J\) through collisions. They do not
give the \(y^2\) statistic's Lipschitz law. At a real double event, the
upper height on the nonreal side is proportional to square-root reverse
time; \(G(d,y)\) is linear in \(y\) at zero, so its contribution has
the same square-root leading scale. No monotonicity follows. A local
factor \((z-r)^2\) gives zero potential, while
\((z-r)^2+d^2\) gives exactly one upper contribution \(G(X-r,d)\).
These are factor checks, not canonical-control constructions.

\subsection{An exact time-zero critical-line and prime decomposition}
The arithmetic reduction in this subsection is specific to \(\lambda=0\)
\cite{MetricComparison,PrimeBaselineAssessment}. Set
\begin{align}\label{eq:green-prime-definitions}
 \mathcal L(X)&=\int_\R K(x-X)
                   \log|\zeta(\tfrac12-\ii x/2)|\dd x,\nonumber\\
 P(X)&=\sum_{n\ge2}\frac{\Lambda_{\rm vM}(n)}{\log n}
          \frac{n^{-5/2}}{1+n^{-4}}\cos\bigl((X/2)\log n\bigr).
\end{align}
Here \(\Lambda_{\rm vM}\) is von Mangoldt, distinct from the Newman
threshold. The logarithmic average is absolutely integrable and the
prime-power series converges absolutely and uniformly for all real \(X\).
\begin{proposition}[Prime baseline with the pole paid]\label{prop:green-prime}
With the canonical normalization,
\begin{equation}\label{eq:green-prime-identity}
 J_0(X)=\mathcal L(X)-P(X)+G(X,1),\qquad
 G(X,1)=\operatorname{atanh}\frac{\sin(\pi/8)}{\cosh(\pi X/8)}.
\end{equation}
For integer \(N\ge2\), if \(P_N\) truncates at \(n\le N\),
\begin{equation}\label{eq:green-prime-tail}
 |P(X)-P_N(X)|\le\tfrac23N^{-3/2}\qquad\text{for every real \(X\)}.
\end{equation}
\end{proposition}
\begin{proof}
Xi reflection permits \(s=1/2-\ii z/2\) in \(H_0(z)=\xi(s)/8\).
Heights \(0,4,8\) become real parts \(1/2,5/2,9/2\). Use
\[
 \xi(s)=\tfrac12s(s-1)\pi^{-s/2}\Gamma(s/2)\zeta(s).
\]
The selected classical identities and gamma growth are recorded in
\cite{DLMFGreenPrime}. The factors \(s\), \(\pi^{-s/2}\),
\(\Gamma(s/2)\), and constants have harmonic logarithmic moduli on
a neighborhood of the closed strip. The zero of \(s\) is at \(z=-\ii\),
outside; gamma is zero-free and pole-free here. Stirling growth on
\(1/4\le\Re(s/2)\le9/4\) is at most linear plus logarithmic for
the log modulus. The exponential Poisson kernels pay this growth and
the vertical-side limit, so these factors have zero
boundary-minus-interior functional. The factor \(s-1\), however, has
a zero at \(z=\ii\) and contributes \emph{\(+G(X,1)\)}. It pays the
zeta pole at \(s=1\); \(\log|\zeta|\) must not be treated as harmonic
across it. Subtracting the paid elementary and gamma logarithms from
\eqref{eq:green-unit-log} proves absolute weighted integrability of
the critical-line logarithm, including real multiple zeros.

On each safe line the absolutely convergent Euler logarithm is
\[
 \log|\zeta(\sigma-\ii x/2)|
   =\sum_{n\ge2}\frac{\Lambda_{\rm vM}(n)}{\log n}
                     n^{-\sigma}\cos\bigl((x/2)\log n\bigr),\quad\sigma>1.
\]
For completeness the substitution \(u=\ee^{2ax}\) and the Euler beta
integral give
\[
 \widehat K(\omega)=\frac1{16a}\int_0^\infty
       \frac{u^{-1/2-\ii\omega/(2a)}}{1+u}\dd u
   =\frac\pi{16a}\operatorname{sech}\frac{\pi\omega}{2a}
   =\tfrac12\operatorname{sech}(4\omega).
\]
At \(\omega=(\log n)/2\) this is \((n^2+n^{-2})^{-1}\).
Consequently the safe-boundary minus interior coefficient is exactly
\begin{equation}\label{eq:green-safe-coefficient}
 \frac{n^{-9/2}}{n^2+n^{-2}}-n^{-5/2}
                         =-\frac{n^{-5/2}}{1+n^{-4}}.
\end{equation}
Evenness of \(K\) retains the displayed cosine with no extra phase or
factor of two. Absolute convergence justifies each interchange and gives
\eqref{eq:green-prime-identity}. Finally
\(\Lambda_{\rm vM}(p^k)/\log(p^k)=1/k\), and the coefficient is zero
elsewhere, so it lies in \([0,1]\). For integer \(N\ge2\),
\[
 |P-P_N|\le\sum_{n>N}n^{-5/2}
       \le\int_N^\infty u^{-5/2}\dd u=\tfrac23N^{-3/2}.
\]
This proves the uniform tail. A real cutoff uses its integer part.
\end{proof}

Positivity gives the unconditional \emph{lower} comparison
\begin{equation}\label{eq:green-prime-lower}
 \mathcal L(X)\ge P(X)-G(X,1).
\end{equation}
To exclude the class in \eqref{eq:green-defect-margin} instead requires
the unproved strict \emph{upper} comparison
\begin{equation}\label{eq:green-missing-comparison}
 \mathcal L(X)<P(X)-G(X,1)+\frac{\sin(a)d}{\cosh(ar)}.
\end{equation}
The prime baseline has a fully paid index tail; the critical-line
logarithmic mean remains the difficult term. The prime tail and the
boundary tail pay separate errors and must fit within the same declared
defect margin, with the height factor retained. No Euler product is
assigned to \(\lambda\ne0\), and no derivative-order or inexpensive
acquisition claim follows from \eqref{eq:green-prime-tail}.

This interface avoids the Gaussian event-amplitude normalization, while
its square-root collision scale remains different from the transverse
\(y^2\) moment. The two missing signed estimates
\eqref{eq:gaussian-missing-comparison} and
\eqref{eq:green-missing-comparison} state the actual remaining
dependencies. Their absence is not a universal exclusion of localized
metrics or source-sensitive arithmetic estimates.
